%%
%% Copyright (C) 2014 Jeffrey P. Hafner <jphafner@buffalo.edu>
%%  
%%  This program is free software: you can redistribute it and/or modify
%%  it under the terms of the GNU General Public License as published by
%%  the Free Software Foundation, either version 3 of the License, or
%%  (at your option) any later version.
%%  
%%  This program is distributed in the hope that it will be useful,
%%  but WITHOUT ANY WARRANTY; without even the implied warranty of
%%  MERCHANTABILITY or FITNESS FOR A PARTICULAR PURPOSE.  See the
%%  GNU General Public License for more details.
%%  
%%  You should have received a copy of the GNU General Public License
%%  along with this program.  If not, see <http://www.gnu.org/licenses/>.
%%


\documentclass[
    11pt,
]{scrartcl}

\usepackage{AMCfull}
\newlength{\mylen}
\usepackage{fp}
\usepackage{amsfonts}
\usepackage{amssymb}

%\graphicspath{{./qbank/collegeBoard/gfx/}}
\graphicspath{{./qbank/nysed/gfx/}}

\AMCnumericOpts{
    hspace=1ex,vspace=1ex,
    vertical=true,reverse=false,
}

%% Begin Document
%%------------------------------------------------
\begin{document}

%\def\AMCntextGoto{\par{\bfseries\emph{Please code the answer on the separate answer sheet.}}}

\element{animals}{
\begin{figure}[p]
  \centering
  %\includegraphics[width=.6\linewidth]{tiger}
  \caption{An animal}
  \AMClabel{tiger}
\end{figure}
\begin{question}{tiger}
    Which is the animal on figure~\AMCref{tiger} page~\AMCpageref{tiger}?
  \begin{choices}
  \correctchoice{A tiger}
    \wrongchoice{A giraffe}
    \wrongchoice{An elephant}
    \wrongchoice{A cat}
  \end{choices}
\end{question}
}

\element{animals}{
\begin{questionmultx}{sqrt}
    \FPeval\VQa{trunc(5+random*15,0)}
    \FPeval\VQs{VQa^0.5}
    Compute $\sqrt{\VQa}$ and round it with two digits after period.
    \AMCnumericChoices{\VQs}{
        digits=3,decimals=2,sign=true,
        borderwidth=0pt,backgroundcol=lightgray,approx=5
    }
\end{questionmultx}
}


\element{animals}{
\begin{question}{inverse}
    \pgfmathrandominteger{\x}{1}{50}
    How much is the reciprocal of $x=\x$?
    \begin{choices}
      \correctchoice{\pgfmathparse{1/\x}\pgfmathresult}
        \wrongchoice{\pgfmathparse{1/(\x +1))}\pgfmathresult}
        \wrongchoice{\pgfmathparse{cos(\x)} \pgfmathresult}
        \wrongchoice{\pgfmathparse{\x^(-0.5)}\pgfmathresult}
    \end{choices}
\end{question}
}

\element{animals}{
\begin{question}{inverse3}
    \pgfmathrandominteger{\x}{1}{50}
    \pgfkeys{/pgf/number format/.cd,fixed,fixed zerofill,precision=3,use comma}
    How much is the reciprocal of $x=\pgfmathprintnumber{\x}$?
    \begin{choices}
        \correctchoice{\pgfmathparse{1/\x}\pgfmathprintnumber{\pgfmathresult}}
        \wrongchoice{\pgfmathparse{1/(\x +1))} \pgfmathprintnumber{\pgfmathresult}}
        \wrongchoice{\pgfmathparse{cos(\x)} \pgfmathprintnumber{\pgfmathresult}}
        \wrongchoice{\pgfmathparse{\x^(-0.5)} \pgfmathprintnumber{\pgfmathresult}}
    \end{choices}
\end{question}
}

\element{animals}{
\begin{questionmult}{graph}
  Let us consider the three functions which graphs are plotted below:
  \pgfmathrandominteger{\a}{2}{4}
  \begin{center}
  \begin{tikzpicture}[domain=0:4]
    \draw[very thin,color=gray] (-0.1,-4.1) grid (3.9,3.9);
    \draw[->] (-0.2,0) -- (4.2,0) node[right] {$x$};
    \draw[->] (0,-4.2) -- (0,4.2) node[above] {$f(x)$};
    \draw[color=red] plot (\x,{(1+\a/4)*\x-\a}) node[right] {$f_{1} (x)$};
    \draw[color=blue] plot (\x,{\a*sin(\x r)}) node[right] {$f_{2}(x)$};
    \draw[color=orange] plot (\x,{\a*cos(\x r)}) node[right] {$f_{3}(x)$};
  \end{tikzpicture}
  \end{center}
  Then:
  \begin{choices}
    \pgfmathrandominteger{\x0}{2}{4}
    \correctchoice{$f_{2}(\x0)$=\pgfmathparse{\a*sin(\x0 r)} \pgfmathprintnumber{\pgfmathresult}.}
    \correctchoice{$f_{3}(\x0)$=\pgfmathparse{\a*cos(\x0 r)} \pgfmathprintnumber{\pgfmathresult}.}
    \wrongchoice{La fonction $f_{1}(x)$ est une fonction linéaire.}
  \end{choices}
\end{questionmult}
}

\element{animals}{
\directlua{math.randomseed (2048)}
\directlua{a=math.random()}
\begin{question}{square-root}
  How much is the square root of \directlua{tex.print(a)}?
  \begin{choices}
  \correctchoice{\directlua{tex.print(math.sqrt(a))}}
    \wrongchoice{\directlua{tex.print(math.sqrt(2*a))}}
    \wrongchoice{\directlua{tex.print(math.sqrt(a*1.001))}}
  \end{choices}
\end{question}
}


%% Begin Quiz
%%------------------------------------------------
\begin{examcopy}[1]

%
\begin{comment}
\begin{minipage}{0.66\linewidth}
    \begin{tabu}{lcl}
        \multicolumn{3}{c}{\LARGE International System of Units (SI)} \\
        \toprule
        Name        & Symbol            & Description \\
        \midrule
        ampere      & \si{\ampere}      & electric current \\
        becquerel   & \si{\becquerel}   & radioactivity \\
        candela     & \si{\candela}     & luminous intensity  \\
        degree Celsius & \si{\degreeCelsius} & temperature relative to \SI{273.15}{\kelvin} \\
        coulomb     & \si{\coulomb}     & electric charge \\
        farad       & \si{\farad}       & electric capacitance \\
        gray        & \si{\gray}        & absorbed dose \\
        henry       & \si{\henry}       & inductance \\
        hertz       & \si{\hertz}       & frequency \\
        joule       & \si{\joule}       & energy, work, heat \\
        katal       & \si{\katal}       & catalytic activity \\
        kilogram    & \si{\kilo\gram}   & mass \\
        kelvin      & \si{\kelvin}      & thermodynamic temperature \\
        lumen       & \si{\lumen}       & luminous flux \\
        lux         & \si{\lux}         & illuminance \\
        meter       & \si{\meter}       & Length \\
        mole        & \si{\mole}        & amount of substance  \\
        newton      & \si{\newton}      & force, weight \\
        ohm         & \si{\ohm}         & electric resistance \\
        pascal      & \si{\pascal}      & pressure, stress \\
        radian      & \si{\radian}      & angle \\
        second      & \si{\second}      & time \\
        siemens     & \si{\siemens}     & electrical conductance \\
        sievert     & \si{\sievert}     & equivalent dose \\
        steradian   & \si{\steradian}   & solid angle \\
        tesla       & \si{\tesla}       & magnetic field strength \\
        volt        & \si{\volt}        & electric potential difference \\
        watt        & \si{\watt}        & power, radiant flux\\
        weber       & \si{\weber}       & magnetic flux \\
        \bottomrule
    \end{tabu}
\end{minipage}
\end{comment}

\begin{minipage}{0.33\linewidth}
    \begin{center}
        \begin{tabu}{lcl}
            \multicolumn{3}{c}{\LARGE Non SI Units} \\
            \toprule
            Name    & Symbol & Definition \\
            \midrule
            minute  & \si{\minute}  & \SI{60}{\second} \\
            hour    & \si{\hour}    & \SI{60}{\minute} \\
            degree  & \si{\degree}  & \SI[fraction-function=\tfrac]{\pi/180}{\radian} \\
            liter   & \si{\liter}   & \SI{1.0}{\deci\meter\cubed} \\
            \bottomrule
        \end{tabu}
    \end{center}
\end{minipage}
\begin{minipage}{0.33\linewidth}
    \begin{center}
        \vskip0.5\baselineskip
        \begin{tabu}{lcl}
            \multicolumn{3}{c}{\LARGE SI Prefixes} \\
            \toprule
            Prefix      & Symbol        & Notation \\
            \midrule
            tera        & \si{\tera}    & \num{e12} \\
            giga        & \si{\giga}    & \num{e9} \\
            mega        & \si{\mega}    & \num{e6} \\
            kilo        & \si{\kilo}    & \num{e3} \\
            hecto       & \si{\hecto}   & \num{e2} \\
            deca        & \si{\deca}    & \num{e1} \\
            deci        & \si{\deci}    & \num{e-1} \\
            centi       & \si{\centi}   & \num{e-2} \\
            milli       & \si{\milli}   & \num{e-3} \\
            micro       & \si{\micro}   & \num{e-6} \\
            nano        & \si{\nano}    & \num{e-9} \\
            pico        & \si{\pico}    & \num{e-12} \\
            \bottomrule
        \end{tabu}
    \end{center}
\end{minipage}
\begin{minipage}{0.33\linewidth}
    \begin{center}
        \begin{tabu}{ll}
            \multicolumn{2}{c}{\LARGE Customary Units} \\
            \toprule
            Name            & Definition \\
            \midrule
            inch    & $\SI{2.54}{\centi\meter}$ \\
            foot    & $\SI{12}{inches}$ \\
            yard    & $\SI{3}{feet}$ \\
            mile    & $\SI{5280}{feet}$ \\
            lb-mass     & $\SI{0.45359237}{\kilo\gram}$ \\
            lb-force    & $\SI{4.4482216152605}{\newton}$ \\
            gallon  & $\SI{3.785411784}{\liter}$ \\
            \bottomrule
        \end{tabu}
    \end{center}
\end{minipage}

\vskip0.5\baselineskip
\begin{minipage}{1.00\linewidth}
    \renewcommand{\arraystretch}{1.25}
    \begin{center}
        \begin{tabu}{X[l]X[l]}
            \multicolumn{2}{c}{\LARGE Kinematic Equations} \\
            \toprule
            Equations       & Definitions \\
            \midrule
            $\bar{v} = \frac{d}{t}$         & $a= \text{acceleration}$ \\
            $a = \frac{\Delta{ }v}{t}$      & $v= \text{velocity}$ \\
            $v_f = v_i + a t$               & $d= \text{displacement or distance}$ \\
            $d = v_i t + \frac{1}{2}at^2$   & $\Delta{}= \text{change}$ \\
            $v^2_f = v^2_i + 2 a d$         & $\bar{v}= \text{average velocity}$ \\
            \bottomrule
        \end{tabu}
    \end{center}
\end{minipage}
\vskip0.5\baselineskip
\begin{minipage}{1.00\linewidth}
    \renewcommand{\arraystretch}{1.0}
    \begin{center}
        \begin{tabu}{X[3l]X[2l]}
            \multicolumn{2}{c}{\LARGE Planetary Measurements} \\
            \toprule
            Name & Value \\
            \midrule
            %Speed of sound in air at STP        & \SI{3.31e2}{\meter\per\second} \\
            mass of Sun                                 & \SI{1.99e30}{\kilo\gram} \\
            mass of Earth                               & \SI{5.98e24}{\kilo\gram} \\
            mass of the Moon                            & \SI{7.35e22}{\kilo\gram} \\
            mean radius of Sun                          & \SI{6.96e8}{\meter} \\
            mean radius of Earth                        & \SI{6.37e6}{\meter} \\
            mean radius of Moon                         & \SI{1.74e6}{\meter} \\
            mean orbital radius of Moon around Earth    & \SI{3.85e8}{\meter} \\
            mean orbital radius of Earth around Sun     & \SI{1.50e11}{\meter} \\
            \bottomrule
        \end{tabu}
    \end{center}
\end{minipage}
\vskip0.5\baselineskip
\begin{minipage}{1.00\linewidth}
    \renewcommand{\arraystretch}{1.0}
    \begin{center}
        \begin{tabu}{X[l]X[l]}
            \multicolumn{2}{c}{\LARGE Selected NIST Physical Constants} \\
            \toprule
            Name & Value \\
            \midrule
            %Boltzmann constant                          & \SI{1.380 648 8 (13) e-23}{\joule\per\kelvin} \\
            %elementary charge                           & \SI{1.602 176 565 (35)e-19}{\coulomb} \\
            %electric constant (exact)                   & \SI{8.854 187 817 e-12}{\farad\per\meter} \\
            %electron mass                               & \SI{9.109 382 91 (40) e-31}{\kilo\gram} \\
            %magnetic constant (exact)                   & \SI{4 \pi e-7}{\henry\per\meter} \\
            %neutron mass                                & \SI{1.674 927 351 (74) e-27}{\kilo\gram} \\
            %Newtonian constant of gravitation           & \SI{6.673 84 (80) e-11}{\meter\cubed\per\kilo\gram\per\second\squared} \\
            %proton mass                                 & \SI{1.672 621 777 (74) e-27}{\kilo\gram} \\
            %Rydberg constant                            & \SI{1.097 373 156 853 9(55) e7}{\per\meter} \\
            %Planck constant                             & \SI{6.626 069 57 (29) e-34}{\joule\second} \\
            speed of light in vacuum (exact)            & \SI{299 792 458}{\meter\per\second} \\
            standard acceleration of gravity (exact)    & \SI{9.806 65}{\meter\per\second\squared} \\
            %standard atmosphere                 & \SI{101 325}{\pascal} \\
            %standard-state pressure             & \SI{100 000}{\pascal} \\
            \bottomrule
        \end{tabu}
    \end{center}
\end{minipage}


\pagebreak


\begin{center}
    Physics at Mergenthaler Topic Assessment
\end{center}
\noindent
\begin{minipage}{10.2em}
    \begin{center}
        Student ID Number:
    \end{center}
    \vspace{-0.50em}
    \AMCcode{StudentNum}{7}
\end{minipage}
\setlength{\mylen}{\textwidth}
\addtolength{\mylen}{-10.5em}
\begin{minipage}{\mylen}
    \begin{center}
        \namefield{\fbox{
            \begin{minipage}{0.98\mylen}
                Last and First Name:
                \vspace*{7mm}\dotfill
                \vspace*{1mm}
            \end{minipage}
        }}
    \end{center}
    \begin{multicols}{2}
        \begin{center}
            \bfseries Instructions:
        \end{center}
        \vspace{-0.5\baselineskip}
        Please use only a pencil on this form.
        Do not make stray marks on this form.
        Print your last and first name in the box provided above.
        Enter your Mergenthaler student ID number by filling in
            the appropriate ovals to the left.
        Enter your chosen selection for each question on this
            form by filling in your selected oval beneath each question.
        %\indent{}\emph{Please note:} This form has a front and back side.
        Questions using the sign \multiSymbole{} may have
            zero, one or several correct answers.
        Other questions have a single correct answer.
        %You will be provided with fifteen ($15$) minutes to complete.
        Correct answers will receive \num{5} points,
            while incorrect answers will receive \num{-1} points.
        It is in your best interest to not guess on any of the questions.
        \emph{Good Luck!}
        %\begin{center}
        %    \bfseries Definitions of select US Customary Units
        %\end{center}
        %\vspace{-1.00\baselineskip}
        %\begin{align*}
        %        \SI{1}{lb} &\equiv \SI{0.45359237}{\kilo\gram} \\
        %        \SI{1}{mile} &\equiv \SI{5280}{feet} \\
        %        \SI{1}{inch} &\equiv \SI{2.54}{\centi\meter} \\
        %        \SI{1}{gallon} &\equiv \SI{3.785 411 784}{\liter} \\
        %\end{align*}
    \end{multicols}
\end{minipage}

\vspace{0.5\baselineskip}
\hrule
\vspace{0.5\baselineskip}

\begin{multicols}{2}
    \insertgroup{animals}

    \begin{question}{inverse}
      \pgfmathrandominteger{\x}{1}{50}
           
      How much is the reciprocal of $x=\x$?

      \begin{choices}
        \correctchoice{\pgfmathparse{1/\x}\pgfmathresult }
        \wrongchoice{\pgfmathparse{1/(\x +1))}\pgfmathresult}
        \wrongchoice{\pgfmathparse{cos(\x)} \pgfmathresult}
        \wrongchoice{\pgfmathparse{\x^(-0.5)}\pgfmathresult}
      \end{choices}
    \end{question}

    \begin{questionmultx}{sqrt}
        \FPeval\VQa{trunc(5+random*15,0)}
        \FPeval\VQs{VQa^0.5}
        %\AMCnumericChoices{\VQs}{digits=3,decimals=2,sign=true,
        %               borderwidth=0pt,backgroundcol=lightgray,approx=5}
        %\AMCnumericChoices{\VQs}{digits=5,decimals=2,sign=true,
        %             borderwidth=0pt,backgroundcol=white,approx=5,vertical=true}
       %\AMCnumericChoices{\VQs}{sign=true,digits=5,decimals=2,scoreexact=3,base=10}
        \AMCnumericChoices{\VQs}{sign=true,digits=5,decimals=2,scoreexact=3,base=10,vertical=false}
      \end{questionmultx}

    \directlua{b=500}
    \directlua{c=110}
    \directlua{d=75}
    \directlua{e=(c*b)/(d+c)}

    \begin{questionmultx}{this}
        Two cities are \SI{\directlua{tex.print(b)}}{\kilo\meter} apart.
        Car A starts at City A and heads towards City B at a constant
            speed of \SI{\directlua{tex.print(c)}}{\kilo\meter\per\hour}.
        Car B starts at City B and heads towards City A at a constant
            speed of \SI{\directlua{tex.print(d)}}{\kilo\meter\per\hour}.
        How far are the two cars from City A when they pass each other?
        %\directlua{tex.print(e)}

        \begin{quote}
        Please enter the significand, $a$, and exponent, $b$,
            of your answer in normalized scientific form,
            $a \times 10^b$, in the provided space.
        \end{quote}

        \AMCnumericChoices{\directlua{tex.print(e)}}{
            vertical=false,
            digits=3,decimals=2,sign=true,Tsign=\hspace{1ex}\Large a,
            borderwidth=0pt,backgroundcol=white,approx=5}
        \AMCnumericChoices{\directlua{tex.print(e)}}{
            vertical=false,
            digits=2,decimals=0,sign=true,Tsign=\hspace{1ex}\Large b,
            borderwidth=0pt,backgroundcol=white,approx=5}
    \end{questionmultx}

    \newcommand\onegrid{
        \begin{tikzpicture}
            %\draw[step=0.5cm] (-3,-3) grid (5,5);
            \draw[step=0.5cm] (0,0) grid (8,8);
        \end{tikzpicture}
    }
    \begin{questionmultx}{drawing}
       Draw something on this grid.
       This is a lot of text.
       This is a lot of text.
       This is a lot of text.
       This is a lot of text.
       This is a lot of text.
       \AMCOpen{contentcommand=onegrid}{\wrongchoice[F]{f}\scoring{0}\correctchoice[J]{j}\scoring{2}}
    \end{questionmultx}

    \directlua{math.randomseed (2048)}
    \directlua{a=math.random()}

    \begin{question}{square-root}
      How much is the square root of \directlua{tex.print(a)}?
      \begin{choices}
      \correctchoice{\directlua{tex.print(math.sqrt(a))}} 
        \wrongchoice{\directlua{tex.print(math.sqrt(2*a))}} 
        \wrongchoice{\directlua{tex.print(math.sqrt(a*1.001))}} 
      \end{choices}
    \end{question}

    \begin{question}{inf-expo-indep}
      \FPeval\VQa{trunc(2 + random * 4,0)}
      \FPeval\VQb{trunc(6 + random * 5,0)}
      \FPeval\VQr{VQa/(VQa+VQb)}
      Let $X$ and $Y$ be two independent random variables,
         following the exponential distribution with
         respective parameters \VQa{} et \VQb{}.
      To which interval does the probability $\mathbb{P}[X<Y]$ belong?
      \begin{multicols}{3}
        \begin{choices}[o]
          \AMCIntervals{\VQr}{0}{1}{0.3}
          %% 1. The correct answer,
          %% 2. The left point of the first interval,
          %% 3. The right point of the last interval,
          %% 4. The length of each interval.
        \end{choices}
      \end{multicols}
    \end{question}

    \begin{questionmult}{graph}
      Let us consider the three functions which graphs are plotted below:
      \pgfmathrandominteger{\a}{2}{4}  
      \begin{center}
        \begin{tikzpicture}[domain=0:4]
          \draw[very thin,color=gray] (-0.1,-4.1) grid (3.9,3.9);
          \draw[->] (-0.2,0) -- (4.2,0) node[right] {$x$};
          \draw[->] (0,-4.2) -- (0,4.2) node[above] {$f(x)$};
          \draw[color=red] plot (\x,{(1+\a/4)*\x-\a}) node[right] {$f_{1} (x)$};
          \draw[color=blue] plot (\x,{\a*sin(\x r)}) node[right] {$f_{2}(x)$};
          \draw[color=orange] plot (\x,{\a*cos(\x r)}) node[right] {$f_{3}(x)$};
        \end{tikzpicture}
      \end{center}
      Then:
      \begin{choices}
        \pgfmathrandominteger{\x0}{2}{4}  
        \correctchoice{$f_{2}(\x0)$=\pgfmathparse{\a*sin(\x0 r)} \pgfmathprintnumber{\pgfmathresult}.}
        \correctchoice{$f_{3}(\x0)$=\pgfmathparse{\a*cos(\x0 r)} \pgfmathprintnumber{\pgfmathresult}.}
        \wrongchoice{La fonction $f_{1}(x)$ est une fonction linéaire.}
      \end{choices}
    \end{questionmult}
\end{multicols}

\clearpage
%%separateanswersheet
\AMCformBegin
\AMCform
\end{examcopy}
\end{document}

